\section{Integer signatures and the softened count}
\label{sec:integer-signatures}

We will use the two-view encoding for all retained cells to obtain
two properties of the softened count.
First, after summing the large cells, the count weight $h_j$ must satisfy
quadratic inequalities that will pass to the binary encoding.  Second,
a tree-value marginal in $C(H)$ must be log-concave, so that attainable
endpoints control the near-boundary violations left in
Proposition~\ref{prop:contamination}.  Both follow from the same integer
balanced-summation argument.

The one-positive-eigenvalue tests are closely related to the
Lorentzian-polynomial framework of Br\"and\'en and Huh~\cite{BH2020}.
The laminar decomposition and balanced-summation argument for widths
at least two appear in~\cite[Lemmas~3.2--3.3 and Proposition~3.4]{OpenAIExactTables2026}.
We prove the closure statement in the form needed here, including
pairs of width one.

\subsection{A projection fact for quadratic forms}

We first record an elementary fact that will be used repeatedly.
If a real symmetric matrix \(Q\) has at most one positive eigenvalue
and \(u^{\mathsf T}Qu>0\), then \(Q\) is nonpositive on
\[
 \{v:v^{\mathsf T}Qu=0\}.
\]
Indeed, a positive-square vector in this subspace, together with \(u\),
would span a two-dimensional positive-definite subspace.
Subtracting the projection of an arbitrary vector \(v\) onto \(u\)
therefore gives
\begin{equation}
 v^{\mathsf T}Qv
 \le
 \frac{(v^{\mathsf T}Qu)^2}{u^{\mathsf T}Qu}.
 \label{eq:signature-projection}
\end{equation}
Conversely, a symmetric form that is nonpositive on a subspace of
codimension at most one has at most one positive eigenvalue.
The latter property is preserved by restriction, by congruence, and
by entrywise limits of matrices of a fixed size.
For limits, this follows from continuity of the ordered eigenvalues.

We also need the boundary case of the projection inequality.  If
$e^{\mathsf T}Qe=0$ but $Qe\ne0$, then $Q$ is nonpositive on
$\{v:v^{\mathsf T}Qe=0\}$.  Indeed, choose $w=Qe$, so
$e^{\mathsf T}Qw=\|Qe\|^2>0$.  For sufficiently small $t>0$,
$u=e+tw$ has positive square.  For $v^{\mathsf T}Qe=0$,
\eqref{eq:signature-projection} gives
\[
 v^{\mathsf T}Qv
 \le \frac{t^2(v^{\mathsf T}Qw)^2}
 {2t e^{\mathsf T}Qw+t^2w^{\mathsf T}Qw}\longrightarrow0.
\]

\subsection{Integer coordinates and balanced summation}

Let \(I\) be a finite coordinate set, let \(m_i\) be nonnegative
integers, and let \(\mathcal L\) be a laminar family of subsets of \(I\):
any two members are nested or disjoint.
For \(G\in\mathcal L\), let
\(\psi_G:\mathbb Z\longrightarrow(0,\infty)\) be log-concave, meaning
\[
 \psi_G(a)^2\ge \psi_G(a-1)\psi_G(a+1)
 \qquad(a\in\mathbb Z).
\]
Fix an integer \(R\).  On the total slice of size \(R\), define
\begin{equation}
 F(z)=
 \prod_{G\in\mathcal L}
 \psi_G\left(\sum_{i\in G}z_i\right)
 \quad\text{if }0\le z_i\le m_i\text{ and }\sum_i z_i=R,
 \label{eq:laminar-weight}
\end{equation}
and set \(F(z)=0\) otherwise.
All coordinates and displays in this subsection are integral.
Coincident members of a laminar family may be combined by multiplying
their factors.

Choose disjoint coordinate pairs
\(\{a_t,b_t\}\), \(t\in J\), with
\(m_{a_t}=m_{b_t}=M_t\).
Summing such a pair \emph{at balance} means summing its profiles
\((k,M_t-k)\), \(0\le k\le M_t\).
Let \(\Phi\) be the weight on the remaining coordinates after summing
all pairs in \(J\) at balance.
Its required total is
\[
 r=R-\sum_{t\in J}M_t.
\]

\begin{lemma}[Integer signature under balanced summation]
\label{lem:integer-signature}
Suppose \(M_t\ge1\) for every pair that is summed.
For every box-valid display \(z\) on the remaining coordinates with
\(\sum_i z_i=r+2\), the symmetric matrix
\begin{equation}
 \bigl(\Phi(z-e_i-e_h)\bigr)_{i,h}
 \quad\text{has at most one positive eigenvalue}.
 \label{eq:integer-removal-signature}
\end{equation}
The diagonal entries in this matrix are the weights after two
removals from the same coordinate; invalid removals have weight zero.
The corresponding assertion for two additions holds for every
box-valid display of total \(r-2\).
\end{lemma}

\begin{proof}
We first prove the removal assertion without summation.
Coordinates with \(z_i=0\) give zero rows and columns and may be
discarded.
Temporarily evaluate all remaining entries using the product of
factors, even when a repeated removal would give a negative coordinate.
This evaluation is well-defined because the factors are defined on
all integers.
Write
\[
 l_G=\sum_{i\in G}z_i,
 \qquad
 \gamma_G=
 \frac{\psi_G(l_G)\psi_G(l_G-2)}
      {\psi_G(l_G-1)^2}.
\]
Log-concavity gives \(0<\gamma_G\le1\).
Factor out the value of the product at the display and, from each
matrix index \(i\), the positive single-removal ratio
\[
 \prod_{G\ni i}\frac{\psi_G(l_G-1)}{\psi_G(l_G)}.
\]
Thus, up to a positive scalar and a positive diagonal congruence,
the entry at \(i,h\) is
\(\prod_{G\supseteq\{i,h\}}\gamma_G\).
Restrict the factor sets to the retained coordinates and combine
duplicates if necessary.  Laminarity then gives the matrix identity
\begin{equation}
 \left(\prod_{G\supseteq\{i,h\}}\gamma_G\right)_{i,h}
 =
 \mathbf 1\mathbf 1^{\mathsf T}
 -
 \sum_{G\in\mathcal L}
 (1-\gamma_G)
 \left(\prod_{\substack{G'\in\mathcal L\\G'\supsetneq G}}
              \gamma_{G'}\right)
 \mathbf 1_G\mathbf 1_G^{\mathsf T}.
 \label{eq:laminar-rank-one}
\end{equation}
For each entry the identity telescopes along the chain of sets
containing both indices.
The right-hand side is a rank-one positive semidefinite matrix
minus a positive semidefinite matrix, so it has at most one positive
eigenvalue.
The only invalid removals still present are repeated removals from a
coordinate with displayed value one.
Their positive diagonal entries must be replaced by zero, which
subtracts a nonnegative diagonal matrix.
This cannot increase the number of positive eigenvalues.

We now induct on the number of summed pairs.
Suppose some pairs have already been summed, and let \(\Psi\) denote
the resulting weight before the next pair of capacity \(M\ge1\) is
summed.
Fix a box-valid display \(z\) on the other coordinates, two above
their new required total.
For \(0\le k\le M\), apply the induction hypothesis to the display
\((z,k,M-k)\).
The resulting matrix, with the two pair coordinates listed last, is
\begin{equation}
 N_k=
 \begin{pmatrix}
  V_k&S(k-1)&S(k)\\
  S(k-1)^{\mathsf T}&u_{k-1}&u_k\\
  S(k)^{\mathsf T}&u_k&u_{k+1}
 \end{pmatrix},
 \label{eq:integer-summation-matrix}
\end{equation}
where
\begin{align}
 (V_k)_{ih}&=\Psi(z-e_i-e_h,k,M-k), \notag\\
 S(l)_i&=\Psi(z-e_i,l,M-1-l), \notag\\
 u_k&=\Psi(z,k-1,M-1-k).
 \label{eq:integer-summation-entries}
\end{align}
Every \(N_k\) has at most one positive eigenvalue.
As usual, all out-of-box terms are zero.
In particular, \(S(-1)=S(M)=0\), and \(u_k=0\) outside
\(1\le k\le M-1\).
The matrix after summation is $\sum_{k=0}^M V_k$.  We must find
one common subspace of codimension at most one on which this sum is
nonpositive; simply adding the matrices $N_k$ would not establish
that property.

If $M=1$, put $S=S(0)$.  After deleting their zero pair axis, $N_0$
and $N_1$ have the respective forms
\[
 \begin{pmatrix}V_0&S\\S^{\mathsf T}&0\end{pmatrix},
 \qquad
 \begin{pmatrix}V_1&S\\S^{\mathsf T}&0\end{pmatrix}.
\]
When $S\ne0$, the null-axis projection just proved shows that both
$V_0$ and $V_1$ are nonpositive on $\{a:a^{\mathsf T}S=0\}$.
Their sum therefore has at most one positive eigenvalue.
If any displayed coordinate $z_i$ is positive, then
$S_i=\Psi(z-e_i,0,0)>0$: the profile has the required total and every
previously summed pair supplies a positive balanced term.  Thus
$S=0$ can occur only when $z=0$, in which case all removals are invalid
and $V_0=V_1=0$.  This proves the width-one case.

Assume henceforth that $M\ge2$.
For each interior index \(k\), \(u_k>0\):
the displayed coordinates are box-valid and have the correct total,
and any balanced profile on each previously summed pair supplies
a positive term in the original positive-factor model.

For larger widths the common subspace is determined by the sum of
the one-removal vectors $S(l)$.  Cumulative sums provide auxiliary
coordinates whose contributions will telescope.  Consider an
old-index vector \(a\) satisfying the single linear condition
\[
 a^{\mathsf T}\sum_{l=0}^{M-1}S(l)=0.
\]
Define
\[
 g_k=
 \frac{a^{\mathsf T}\sum_{l=0}^{k-1}S(l)}{u_k}
 \quad(1\le k\le M-1),
 \qquad
 g_k=0\quad\text{otherwise}.
\]
The endpoint condition on \(a\), as well as the cumulative-sum
definition at interior indices, gives
\begin{equation}
 a^{\mathsf T}S(l)=u_{l+1}g_{l+1}-u_lg_l
 \qquad(0\le l\le M-1).
 \label{eq:integer-cumulative}
\end{equation}
For \(v_k=(a,g_k,-g_k)\), its products under the form \(N_k\) with
the first and second new coordinate axes are, respectively,
\[
 u_{k-1}(g_k-g_{k-1}),
 \qquad
 u_{k+1}(g_{k+1}-g_k).
\]
If either of those axes has positive square, applying
\eqref{eq:signature-projection} to that axis and adding the other
nonnegative term gives
\begin{equation}
 v_k^{\mathsf T}N_kv_k
 \le
 u_{k-1}(g_k-g_{k-1})^2+
 u_{k+1}(g_k-g_{k+1})^2.
 \label{eq:integer-projection-bound}
\end{equation}
If both axis squares vanish, their diagonal entries
\(u_{k-1},u_{k+1}\) are zero.
The endpoint cases \(k=0,M\) have a positive other diagonal because
\(M\ge2\); hence in the remaining case \(u_k>0\).
The sum of the two axes then has square \(2u_k>0\) and product zero
with \(v_k\).
Applying \eqref{eq:signature-projection} to this sum again gives
\eqref{eq:integer-projection-bound}, whose right-hand side is now
zero.
This includes the case \(M=2,k=1\).

Substitute \eqref{eq:integer-cumulative} into
\eqref{eq:integer-projection-bound} and expand.  The terms involving
$g_kg_{k-1}$ and $g_kg_{k+1}$ cancel, leaving
\begin{equation}
 a^{\mathsf T}V_ka
 +2u_kg_k^2
 -u_{k-1}g_{k-1}^2
 -u_{k+1}g_{k+1}^2
 \le0.
 \label{eq:integer-telescoping}
\end{equation}
Sum this inequality over \(0\le k\le M\).
Every interior \(u_kg_k^2\) cancels, and all exterior terms are zero.
Consequently \(\sum_kV_k\), which is the desired matrix after summing
the next pair, is nonpositive on the kernel of one linear functional.
Its number of positive eigenvalues is therefore at most one.
This completes the induction.

Finally reflect every coordinate in its box:
\(z_i\mapsto m_i-z_i\).
A factor on a set \(G\) remains a positive log-concave sequence
in the sum on \(G\), with its argument reversed and translated.
A balanced pair of capacity \(M\) remains balanced under this
reflection, so reflection commutes with balanced summation.
The reflected remaining weight is
\(\Phi^\vee(z)=\Phi(m-z)\).
A display two below the original total becomes a display two above
the reflected total, and
\[
 \Phi^\vee(m-z-e_i-e_h)=\Phi(z+e_i+e_h).
\]
The removal assertion for the reflected model is precisely the
addition assertion for the original model.
\end{proof}

The conclusion of Lemma~\ref{lem:integer-signature} also holds for
entrywise limits of these weights.
All sums used in this section are finite, so such limits commute
with the summations.
For a member $G$ of the laminar family, an exact equality
\(\sum_{i\in G}z_i=l\) can be imposed by multiplying its factor by
the positive log-concave sequence
\[
 u^{|\sum_{i\in G}z_i-l|},
 \qquad 0<u<1,
\]
followed by \(u\downarrow0\).
We will take limits only after applying the positive-factor argument,
including all divisions in its proof.

\subsection{Application to the softened table weights}

Recall the weights \(h_j(x,y)\) from \eqref{eq:hj}, on small-cell
counts of total
\[
 p=\sum_{s\ \mathrm{small}}b_s.
\]
Here and below each small cell contributes two count coordinates,
one for its row view and one for its column view.

\begin{corollary}
\label{cor:hj-signatures}
For \(0\le j\le H\), the matrix of two-removal weights of \(h_j\)
at every box-valid display of total \(p+2\) has at most one positive
eigenvalue.
The corresponding two-addition matrix at every box-valid display
of total \(p-2\) has the same property.
\end{corollary}

\begin{proof}
Introduce a finite physical interval
\([\ell_s,\ell_s+M_s]\) for every retained cell.
For small and free cells take \([0,b_s]\), and for a tree cell
take
\[
 [\ell_s,\ell_s+M_s]=[-10db_s,10db_s].
\]
Lemma~\ref{lem:forest} shows that these tree intervals contain
every completion required by the definition of \(h_j\).
Represent a physical row-view value as \(\ell_s+x_s\) and
a physical column-view value as \(\ell_s+M_s-y_s\), where
\(0\le x_s,y_s\le M_s\).
The two views agree when \(x_s+y_s=M_s\).
Use the full total slice
\[
 \sum_s(x_s+y_s)=\sum_sM_s.
\]

For a row or column vertex $v$, let $\kappa_v$ be its prescribed
margin and write its physical line sum as
\[
 L_v(x,y)=
 \begin{cases}
 \sum_{s\ni v}(\ell_s+x_s),&v\text{ is a row},\\
 \sum_{s\ni v}(\ell_s+M_s-y_s),&v\text{ is a column}.
 \end{cases}
\]
For $0<\alpha<1$, give each box-valid profile on the full total
slice the positive weight
\begin{equation}\label{eq:laminar-table-model}
 \alpha^{\sum_{v\text{ nonroot}}|L_v-\kappa_v|}
 2^{-\frac{jA}{H}\left(
   \sum_{v\text{ root}}|L_v-\kappa_v|
   +\sum_{t\in\mathcal T}
      \frac{\dist(\ell_t+x_t,[0,b_t])}{b_t}\right)}.
\end{equation}
Each factor is log-concave in its indicated coordinate sum.
Row factors use disjoint groups of $x$ coordinates, and column factors
use disjoint groups of $y$ coordinates.  These two collections are also
disjoint from one another.  The only additional factors are unary tree
factors, each nested in its row group.  Thus the entire family is
laminar.  Letting $\alpha\downarrow0$ imposes every nonroot equation
exactly while retaining the root and tree penalties.

Sum every large-cell pair at balance.
The widths of all summed pairs are positive.
The remaining total is \(p\).
For every small count profile on that slice and every free-cell
assignment, the exact nonroot equations have a unique tree
completion, and it lies in the chosen intervals by
Lemma~\ref{lem:forest}.
Therefore, in the equality limit, each free assignment contributes
exactly one term, with precisely the root and tree penalty defining
\(h_j\); there is no multiplicity from tree coordinates.
The removal assertion follows from
Lemma~\ref{lem:integer-signature} and passage to a finite limit.
Reflecting all coordinates, including those that are summed, gives
the addition assertion by the same lemma.
\end{proof}

\subsection{Tree marginals and contamination}

For a tree cell \(t\), let \(a_t(k)\) be the unnormalized marginal
mass of its physical value \(k\) under the final weighted sum
\(C(H)\).
Explicitly,
\begin{equation}
 a_t(k)=
 \sum_{\substack{x_s+y_s=b_s\ \text{for every small }s\\
                  0\le z_e\le b_e\ \text{for every free }e\\
                  z_t(x,y,z)=k}}
 2^{-D(x,y,z)}.
 \label{eq:tree-marginal-definition}
\end{equation}
The sums are over integer profiles, and the tree value in the
constraint is the unique nonroot completion.
In particular \(\sum_k a_t(k)=C(H)\).

\begin{lemma}[Log-concave tree marginals]
\label{lem:tree-marginal}
For integers \(v<k<w\) such that \(a_t(v),a_t(w)>0\),
\begin{equation}
 a_t(k)\ge
 a_t(v)^{(w-k)/(w-v)}
 a_t(w)^{(k-v)/(w-v)}.
 \label{eq:tree-marginal-interpolation}
\end{equation}
Consequently the positive support of \(a_t\) is an integer interval.
\end{lemma}

\begin{proof}
Use the finite interval representation and positive weight
\eqref{eq:laminar-table-model} at $j=H$, with $0<\alpha<1$.
Sum every cell pair except $t$ at balance.  All retained cells have
positive width, including small cells of bound one, so
Lemma~\ref{lem:integer-signature} applies.  The remaining required
total is $M_t$.
Let $a_{t;\alpha}(k)$ be the resulting weight of the retained balanced
pair with physical value $k$, for
$\ell_t\le k\le\ell_t+M_t$.
It is positive throughout this integer interval: after fixing the
retained pair, any balanced choices for all other pairs satisfy the
boxes and total slice, and every factor is positive.  Exact nonroot
equations are not required before taking the limit.

At an interior value \(k\), the retained profile is
\[
 (k-\ell_t,\ \ell_t+M_t-k).
\]
Adding one to both coordinates gives a box-valid display two above
its required total.
Lemma~\ref{lem:integer-signature} applied to the two remaining
coordinates gives the matrix
\[
 \begin{pmatrix}
  a_{t;\alpha}(k-1)&a_{t;\alpha}(k)\\
  a_{t;\alpha}(k)&a_{t;\alpha}(k+1)
 \end{pmatrix}
\]
with at most one positive eigenvalue.
Its positive trace implies a nonpositive determinant, so
\[
 a_{t;\alpha}(k)^2
 \ge a_{t;\alpha}(k-1)a_{t;\alpha}(k+1).
\]
All these marginal masses are positive.
Their logarithms therefore have nonincreasing successive
differences, which proves global interpolation between any two
values in the interval.

Now let $\alpha\downarrow0$.  The sums are finite, and the surviving
profiles satisfy every nonroot equation.  For each balanced small
profile and free assignment those equations have exactly one tree
completion, contained in the chosen intervals by
Lemma~\ref{lem:forest}.  The surviving weight is exactly $2^{-D}$.
Consequently, term by term,
\[
 a_{t;\alpha}(k)\longrightarrow a_t(k).
\]
Pass the global interpolation inequality to this limit.
For positive limiting endpoint masses its right-hand side is
positive, giving \eqref{eq:tree-marginal-interpolation} and ruling
out any internal support gap.
No inference from adjacent inequalities alone after the limit is
needed.
\end{proof}

We can now finish the comparison between the softened count and the
number of feasible tables.  The endpoint attainability obtained by
tightening supplies the two positive marginal masses needed below.

\begin{proof}[Proof of Proposition~\ref{prop:contamination}]
Expand $C(H)$ over balanced small profiles and free integer vectors,
giving each assignment mass $2^{-D}$.  The forest equations determine
one tree vector for each such assignment.  It represents a feasible
table exactly when every root deviation is zero and every tree value
belongs to its allowed interval.  These assignments each have mass
one, so $C(H)\ge Z\ge1$.  The total number of assignments is
$\prod_{s\in\mathcal S}(b_s+1)\prod_{e\in\mathcal F}(b_e+1)\le2^K$,
and every mass is at most one.

Call an assignment \emph{far} if a root deviation is nonzero or a tree
value has distance at least $\gamma b_t$ from $[0,b_t]$.  A nonzero
root deviation is an integer of absolute value at least one.
Since $\gamma<1$ and $A\gamma\ge4K$, every far assignment has
$D\ge4K$.  The total mass of all far assignments is therefore at most
$2^K2^{-4K}=2^{-3K}\le\epsilon/128$.  Because $C(H)\ge1$, this also
bounds its fraction of the total mass.

For each tree cell $t$, its marginal satisfies
\[
\sum_{k\in\mathbb Z}a_t(k)=C(H),\qquad
0\le a_t(k)\le2^K,\qquad a_t(0),a_t(b_t)\ge1.
\]
The endpoint inequalities follow from endpoint attainability after
tightening and the unit mass of each feasible table.
If $k<0$ and $a_t(k)>0$, use
\eqref{eq:tree-marginal-interpolation} from $k$ toward $b_t$.
For every integer
$0\le r\le\lfloor(b_t-k)/K\rfloor$, writing
$\alpha=r/(b_t-k)\le1/K$, the endpoint cases and interpolation give
\[
a_t(k+r)\ge a_t(k)^{1-\alpha}a_t(b_t)^\alpha
\ge a_t(k)\,2^{-K\alpha}\ge\frac{a_t(k)}2.
\]
All these integers lie between $k$ and $b_t$.  Their number is
$\lfloor(b_t-k)/K\rfloor+1\ge b_t/K$, and hence
$a_t(k)/C(H)\le2K/b_t$.

For the other exterior side, let $k>b_t$ with $a_t(k)>0$ and
interpolate toward zero.  Every integer
$0\le r\le\lfloor k/K\rfloor$ satisfies
\[
a_t(k-r)\ge a_t(k)^{1-r/k}a_t(0)^{r/k}
\ge a_t(k)\,2^{-Kr/k}\ge\frac{a_t(k)}2.
\]
There are $\lfloor k/K\rfloor+1\ge b_t/K$ such integers, giving
the same bound $a_t(k)/C(H)\le2K/b_t$.  The bound also holds when
$a_t(k)=0$.

On either exterior side of $[0,b_t]$, the number of integers at
positive distance strictly less than $\gamma b_t$ is
$\max\{\lceil\gamma b_t\rceil-1,0\}<\gamma b_t$.
Consequently the fraction of mass with
$0<\dist(z_t,[0,b_t])<\gamma b_t$ is at most $4\gamma K$.
A union bound over at most $d$ tree cells bounds all these near
violations by
$4d\gamma K=\epsilon/128$.  Every infeasible assignment is far or
has such a near tree violation.  We conclude that
\[
\frac{C(H)-Z}{C(H)}\le\frac{\epsilon}{64}.
\]
Finally,
$C(H)\le Z/(1-\epsilon/64)\le(1+\epsilon/32)Z$ for
$0<\epsilon<1$, as claimed.
\end{proof}
